\subsecao{Visão Geral da Avaliação Experimental}

A avaliação do sistema autônomo de captura de armadilhas de vetores biológicos foi realizada por meio de um conjunto integrado de testes práticos, organizado para verificar se os sete objetivos específicos definidos na Introdução foram cumpridos. Os experimentos combinaram avaliações numéricas do desempenho das redes neurais sobre bancos de imagens sintéticas criados para o projeto, testes estatísticos por simulação de Monte Carlo para avaliar o filtro de rastreamento e baterias completas de voos autônomos integrados no simulador físico Gazebo Sim conectado ao piloto automático real ArduPilot no formato \textit{Software-in-the-Loop} (SITL).

Os resultados e as discussões estão organizados em cinco etapas principais:
\begin{enumerate}
	\item Comparação da precisão, da capacidade de adaptação a novas imagens (generalização) e do tempo de resposta (latência) das redes da família YOLO em cada uma das três fases visuais (Objetivos 1 e 2);
	\item Desempenho estatístico e validação em voo simulado do Filtro de Kalman Estendido (EKF) quanto à redução de ruídos e à capacidade de manter a trajetória quando a câmera perde o alvo de vista (Objetivo 3);
	\item Comportamento da margem de segurança na tela (histerese visual) para manter a armadilha sempre visível dentro do campo de visão da câmera sem interrupções (Objetivo 4);
	\item Estabilidade e precisão do alinhamento de giro horizontal (\textit{Yaw}) durante a aproximação final (Objetivo 5);
	\item Consolidação estatística da missão autônoma completa nos ensaios de Monte Carlo, avaliando a precisão em milímetros do encaixe e a taxa global de sucesso das capturas (Objetivos 6 e 7).
\end{enumerate}

\subsecao{Teste 1: Avaliação Comparativa e Desempenho dos Detectores YOLO (Objetivos 1 e 2)}

O primeiro teste avaliou em números quatro arquiteturas leves da família YOLO (YOLOv8n, YOLOv10n, YOLOv11n e YOLO26n) nas três etapas visuais da missão. O objetivo foi verificar se os modelos atingiam a meta de precisão na detecção ($mAP@50 > 90\%$) e a velocidade mínima de processamento em tempo real (acima de $20\text{ fps}$, o que equivale a um tempo de resposta menor que $50\text{ ms}$ por imagem analisada individualmente, com \textit{batch} = 1).

A Tabela~\ref{tab:comparativo_yolo_fase1_resultados} reúne os resultados obtidos pelas redes no conjunto de teste independente da Fase 1 (detecção geral da armadilha em alturas entre 20,0 m e 2,0 m).

\begin{table}[H]
	\centering
	\footnotesize
	\setlength{\tabcolsep}{3.5pt}
	\captionsetup{justification=centering}
	\caption{Comparativo de desempenho entre arquiteturas da família YOLO na detecção macroscópica da armadilha na Fase 1}
	\label{tab:comparativo_yolo_fase1_resultados}
	\vspace{0.2cm}
	\begin{tabular}{lcccccc}
		\hline
		\textbf{Arquitetura} & \textbf{Precisão ($P$)} & \textbf{\textit{Recall} ($R$)} & \textbf{$mAP@50$} & \textbf{$mAP@50\text{--}95$} & \textbf{Latência (ms)} & \textbf{Taxa (fps)} \\
		\hline
		YOLOv8n    & 100,00\%     & 99,96\%    & 99,50\%    & 71,11\%         & 42,6           & 23,5 \\
		YOLOv10n   & 96,01\%      & 97,60\%    & 99,20\%    & 72,54\%         & 35,0           & 28,6 \\
		YOLOv11n   & 99,19\%      & 99,60\%    & 99,47\%    & 73,38\%         & 44,7           & 22,4 \\
		\textbf{YOLO26n} & \textbf{98,39\%} & \textbf{97,77\%} & \textbf{99,23\%} & \textbf{75,68\%} & \textbf{19,3} & \textbf{51,8} \\
		\hline
	\end{tabular}
	\vspace{0.2cm}
	\begin{flushleft}
		\footnotesize Fonte: elaboração própria.
	\end{flushleft}
\end{table}

Na Fase 1, todos os quatro modelos testados alcançaram a meta de $mAP@50 > 90\%$, ficando acima de $99{,}20\%$. Quando avaliados pelo índice mais rigoroso $mAP@50\text{--}95$ (que mede o encaixe exato da caixa de detecção considerando níveis de sobreposição IoU entre 0,50 e 0,95), o modelo YOLO26n teve o melhor resultado ($75{,}68\%$, superando o YOLOv8n em 4,57 pontos percentuais). Quanto à velocidade de processamento, o YOLO26n respondeu em apenas $19{,}3\text{ ms}$ por imagem, o que corresponde a $51{,}8\text{ fps}$, superando com folga o mínimo exigido de 20 fps. Essa rapidez vem da otimização dos blocos de processamento interno da arquitetura YOLO26n, que reduzem o uso de memória sem perder a capacidade de enxergar bem alvos distantes. Por unir alta precisão geométrica em um campo de visão amplo com baixo tempo de processamento, o YOLO26n foi o escolhido para a Fase 1.

A Tabela~\ref{tab:comparativo_yolo_fase2_resultados} resume o desempenho das redes no banco de imagens da Fase 2 (detecção simultânea das alças de içamento e da tampa da armadilha entre 2,0 m e 1,50 m).

\begin{table}[H]
	\centering
	\footnotesize
	\setlength{\tabcolsep}{3.5pt}
	\captionsetup{justification=centering}
	\caption{Comparativo de desempenho entre arquiteturas da família YOLO na detecção multi-classe da Fase 2 (alças e tampa)}
	\label{tab:comparativo_yolo_fase2_resultados}
	\vspace{0.2cm}
	\begin{tabular}{lcccccc}
		\hline
		\textbf{Arquitetura} & \textbf{Precisão ($P$)} & \textbf{\textit{Recall} ($R$)} & \textbf{$mAP@50$} & \textbf{$mAP@50\text{--}95$} & \textbf{Latência (ms)} & \textbf{Taxa (fps)} \\
		\hline
		YOLOv8n    & 99,74\%      & 99,88\%    & 99,50\%    & 87,49\%         & 42,6           & 23,5 \\
		\textbf{YOLOv10n} & \textbf{98,02\%} & \textbf{98,58\%} & \textbf{99,43\%} & \textbf{89,08\%} & \textbf{22,4} & \textbf{44,6} \\
		YOLOv11n   & 98,84\%      & 99,84\%    & 99,49\%    & 88,23\%         & 44,7           & 22,4 \\
		YOLO26n    & 98,69\%      & 98,70\%    & 99,46\%    & 88,80\%         & 19,3           & 51,8 \\
		\hline
	\end{tabular}
	\vspace{0.2cm}
	\begin{flushleft}
		\footnotesize Fonte: elaboração própria.
	\end{flushleft}
\end{table}

Na Fase 2, o principal desafio foi diferenciar três alças muito próximas dispostas em um triângulo assimétrico. O modelo YOLOv10n atingiu o índice $mAP@50\text{--}95$ de $89{,}08\%$ e velocidade de $44{,}6\text{ fps}$ ($22{,}4\text{ ms}$). O motivo técnico para escolher o YOLOv10n nessa fase está na sua estrutura com atribuição dupla de rótulos (\textit{Dual Label Assignment}), que permite rodar de forma nativa sem o filtro tradicional de Supressão de Não-Máximos (livre de NMS). Em detectores comuns que dependem do NMS (como o v8n e o v11n), a proximidade de poucos milímetros entre as alças triangulares vistas de ângulos inclinados pode fazer o sistema apagar uma das caixas vizinhas por confundi-la com uma detecção repetida da mesma peça. Ao eliminar a etapa de NMS durante a análise, o YOLOv10n evitou o desaparecimento intermitente das pontas do triângulo, garantindo um cálculo estável do ângulo de giro do drone a 20 Hz.

Por fim, a Tab.~\ref{tab:comparativo_yolo_etapa3_alinhamento} mostra os resultados do treinamento das quatro arquiteturas no conjunto da Fase 3 (alinhamento fino em altura muito baixa, entre 0,35 m e 0,20 m, focado apenas no aro circular central interno da tampa).

\begin{table}[H]
	\centering
	\footnotesize
	\setlength{\tabcolsep}{3.5pt}
	\captionsetup{justification=centering}
	\caption{Comparativo de desempenho entre arquiteturas YOLO no treinamento da Fase 3 (alinhamento fino a 0,20 m)}
	\label{tab:comparativo_yolo_etapa3_alinhamento}
	\vspace{0.2cm}
	\begin{tabular}{lcccccc}
		\hline
		\textbf{Modelo} & \textbf{Precisão ($P$)} & \textbf{\textit{Recall} ($R$)} & \textbf{$mAP@50$} & \textbf{$mAP@50\text{--}95$} & \textbf{\textit{Box Loss}} & \textbf{Latência (ms)} \\
		\hline
		YOLOv8n   & 99,93\%    & 100,0\%    & 99,50\%    & 95,75\%         & 0,3069             & 42,6 \\
		\textbf{YOLOv10n}  & \textbf{99,93\%} & \textbf{100,0\%} & \textbf{99,50\%} & \textbf{99,50\%} & \textbf{0,2584} & \textbf{35,0} \\
		YOLOv11n  & 99,93\%    & 100,0\%    & 99,50\%    & 99,50\%         & 0,2162             & 44,7 \\
		YOLO26n   & 99,93\%    & 100,0\%    & 99,50\%    & 99,50\%         & 0,2331             & 19,3 \\
		\hline
	\end{tabular}
	\vspace{0.2cm}
	\begin{flushleft}
		\footnotesize Fonte: elaboração própria.
	\end{flushleft}
\end{table}

A criação de um modelo exclusivo para a Fase 3 melhorou ainda mais os índices de validação, com três das quatro redes atingindo $99{,}50\%$ na métrica $mAP@50\text{--}95$. Esse ganho vem diretamente da forma como o problema foi simplificado: ao trocar o corpo externo da armadilha de 255 mm (que acabava cortado nas bordas da câmera \textit{Full HD}) pelo aro circular interno de $72\text{ mm} \times 72\text{ mm}$, reduziram-se os problemas causados pela mudança brusca de tamanho na tela e pelas sombras projetadas pelo corpo do drone. A $0{,}20\text{ m}$ de altura em relação ao alvo, o aro central forma um círculo bem definido no meio da imagem ($5{,}5\text{ px/mm}$), permitindo calcular o centro exato de encaixe com desvio mínimo.

O YOLOv10n foi mantido como o modelo oficial da Fase 3 por reunir o índice $mAP@50\text{--}95$ de $99{,}50\%$, erro muito baixo no ajuste da caixa de detecção ($\text{Box Loss} = 0{,}2584$), tempo de resposta estável de $35{,}0\text{ ms}$ ($28{,}6\text{ fps}$) e ausência total de NMS, o que evita variações no tempo de processamento logo antes da descida final sem câmera.

A Figura~\ref{fig:predicoes_yolo_tres_fases} mostra exemplos reais das detecções feitas pelos modelos escolhidos nas três etapas da missão, destacando como as caixas de marcação se comportam em diferentes alturas de voo.

\begin{figure}[H]
	\centering
	\captionsetup{justification=centering}
	\caption{Exemplos de inferência e caixas delimitadoras preditas pelos detectores selecionados nas três fases}
	\label{fig:predicoes_yolo_tres_fases}
	\vspace{0.2cm}
	\includegraphics[width=0.92\linewidth]{figuras/predicoes_validacao_yolo.png}
	\vspace{0.2cm}
	\begin{flushleft}
		\footnotesize Fonte: elaboração própria.
	\end{flushleft}
\end{figure}

\subsubsecao{Discussão Crítica e Fatores de Melhoria dos Detectores Visuais}

Embora as redes neurais tenham superado todas as metas numéricas no ambiente de simulação física, levar esse sistema de visão computacional para testes reais em campo traz limitações práticas que precisam ser analisadas:
\begin{enumerate}
	
	\item \textbf{Condições de Iluminação Extrema e Sombras Complexas:} No simulador Gazebo, a luz do sol é reproduzida como uma fonte direta única com sombras calculadas matematicamente. Já em ambientes abertos reais, a luz espalhada em dias nublados, o brilho intenso vindo de reflexos nas superfícies plásticas molhadas pela chuva e o movimento da vegetação balançando com o vento criam mudanças de contraste que não foram reproduzidas nas imagens sintéticas. Para superar essa limitação, recomenda-se em trabalhos futuros aplicar técnicas de variação aleatória de domínio (\textit{Domain Randomization}), conforme formulado por Tobin \textit{et al.}~\cite{tobin2017domain}, variando estocasticamente a intensidade da luz, texturas e renderização realista baseada em física (PBR), além da montagem de bancos de imagens híbridos com modelos generativos de difusão. Se viável, ainda é recomendado montar um novo \textit{dataset} com imagens da armadilha real;
	
	\item \textbf{Quantização e Otimização para Hardware Embarcado:} Os tempos de resposta (latências) medidos nos experimentos foram obtidos em uma placa de vídeo (GPU) de desenvolvimento. Para embarcar esse sistema no drone utilizando computadores compactos e de baixo consumo de energia, será importante converter os parâmetros treinados da rede para um formato mais leve com precisão inteira de 8 bits (quantização \textit{INT8}). Estudos recentes da área mostram que a quantização INT8 em modelos YOLO compactos mantém mais de $98\%$ da precisão original enquanto reduz o tempo de processamento pela metade, o que pode garantir taxas acima de 40 fps mesmo quando o equipamento trabalha sob fortes limitações de aquecimento e energia \cite{rey2025performance}.
	
\end{enumerate}

\subsecao{Teste 2: Validação em Malha Fechada do Filtro EKF (Objetivo 3)}

O segundo teste avaliou na prática o desempenho do Filtro de Kalman Estendido (EKF), verificando sua capacidade de reduzir as oscilações rápidas (\textit{jitter}) nas detecções da câmera, suavizar as variações de ângulo na transição circular completa ($S^1$) e manter a estimativa da trajetória do drone durante bloqueios temporários de visão de até 4,0 segundos, como definido no Objetivo 3.

A avaliação começou com uma bateria independente de $N = 50$ simulações estatísticas de Monte Carlo executadas diretamente sobre os módulos reais \textit{FiltroAlvoEKF} e \textit{FiltroAnguloEKF}. Os resultados gerais estão reunidos na Tab.~\ref{tab:validacao_ekf_resultados}.

\begin{table}[H]
	\centering
	\footnotesize
	\setlength{\tabcolsep}{2pt}
	\captionsetup{justification=centering}
	\caption{Métricas estatísticas consolidadas da validação estocástica do Filtro EKF ($N = 50$ ensaios de Monte Carlo)}
	\label{tab:validacao_ekf_resultados}
	\vspace{0.2cm}
	\begin{tabular}{>{\raggedright\arraybackslash}p{3.8cm}cccc}
		\hline
		\textbf{Métrica de Avaliação} & \textbf{Média $\pm$ Desv. Padrão} & \textbf{Faixa [Mín; Máx]} & \textbf{Meta} & \textbf{Resultado} \\
		\hline
		RMSE Posição Bruta ($XY$) & $58{,}57 \pm 2{,}97\text{ mm}$ & $[52{,}88; 68{,}57]\text{ mm}$ & \textit{Baseline} & Referência \\
		RMSE Posição com EKF ($XY$) & $28{,}17 \pm 3{,}57\text{ mm}$ & $[22{,}19; 39{,}29]\text{ mm}$ & $<$ \textit{Baseline} & Atendido \\
		\textbf{Atenuação de \textit{Jitter} em $XY$} & $\mathbf{51{,}98 \pm 5{,}00\%}$ & $[\mathbf{37{,}58; 61{,}28\%}]$ & $\mathbf{\ge 40{,}0\%}$ & \textbf{Superado} \\
		RMSE Guinada Bruto (\textit{Yaw}) & $4{,}00^\circ \pm 0{,}29^\circ$ & $[3{,}28^\circ; 4{,}58^\circ]$ & \textit{Baseline} & Referência \\
		RMSE Guinada com EKF (\textit{Yaw}) & $1{,}91^\circ \pm 0{,}30^\circ$ & $[1{,}37^\circ; 2{,}52^\circ]$ & $<$ \textit{Baseline} & Atendido \\
		\textbf{Atenuação de Ruído de \textit{Yaw}} & $\mathbf{52{,}23 \pm 6{,}34\%}$ & $[\mathbf{37{,}81; 66{,}46\%}]$ & $\mathbf{\ge 40{,}0\%}$ & \textbf{Superado} \\
		Deriva em Oclusão Total de 4,0 s & $127{,}55 \pm 59{,}97\text{ mm}$ & $[34{,}46; 265{,}16]\text{ mm}$ & $\le 250{,}0\text{ mm}$ & Atendido \\
		Deriva de \textit{Yaw} em 4,0 s & $7{,}65^\circ \pm 4{,}90^\circ$ & $[0{,}01^\circ; 20{,}82^\circ]$ & Informativo & Atendido \\
		Erro de Reaquisição pós-Oclusão & $23{,}37 \pm 4{,}60\text{ mm}$ & $[14{,}80; 35{,}87]\text{ mm}$ & $\le 60{,}0\text{ mm}$ & Atendido \\
		Erro de Convergência pós-\textit{Reset} & $19{,}13 \pm 6{,}26\text{ mm}$ & $[9{,}32; 34{,}43]\text{ mm}$ & $\le 80{,}0\text{ mm}$ & Atendido \\
		Sustentação sem Divergência (4 s) & $100{,}0\% \pm 0{,}0\%$ & $[100{,}0; 100{,}0\%]$ & $100{,}0\%$ (80 ciclos) & Atendido \\
		\textit{Reset} Determinístico pós-4,0 s & $100{,}0\% \pm 0{,}0\%$ & $[100{,}0; 100{,}0\%]$ & $100{,}0\%$ dos ensaios & Atendido \\
		\hline
	\end{tabular}
	\vspace{0.2cm}
	\begin{flushleft}
		\footnotesize Fonte: elaboração própria.
	\end{flushleft}
\end{table}

Os resultados confirmam que a redução das oscilações rápidas (\textit{jitter}) chegou a $51{,}98\%$ nos movimentos horizontais e a $52{,}23\%$ no ângulo de giro horizontal (guinada), superando a meta mínima de $40\%$ de melhoria em relação aos dados brutos vindos da rede neural. O erro médio de posição do centro do alvo caiu de $58{,}57\text{ mm}$ para $28{,}17\text{ mm}$, enquanto o ruído na medição do giro diminuiu de $4{,}00^\circ$ para $1{,}91^\circ$. Durante o intervalo de 4,00 segundos sem imagem (80 ciclos seguidos a 20 Hz em que o detector visual foi desligado propositalmente), o cálculo preditivo do EKF manteve o desvio médio acumulado em apenas $127{,}55\text{ mm}$ (bem abaixo do limite de 250 mm). Assim que a câmera voltou a enxergar o alvo, o filtro absorveu as novas leituras e corrigiu o erro de posição para $23{,}37\text{ mm}$ logo no primeiro ciclo, sem causar trancos na trajetória do drone. Nos casos em que o bloqueio visual ultrapassou o limite de 4,0 segundos, o sistema de segurança acionou a reinicialização automática (\textit{reset}) em $100\%$ dos testes, garantindo a proteção contra estimativas sem visão por tempo indeterminado.

Em paralelo, o filtro foi testado em voo simulado completo no Gazebo Sim, aplicando um corte proposital de visão de $3{,}76\text{ segundos}$ durante a descida intermediária (entre 1,29 m e 0,88 m de altura). Os dados gravados pela telemetria mostraram uma redução de oscilações (\textit{jitter}) de $85{,}0\%$ na posição horizontal e de $91{,}2\%$ no ângulo de giro durante a aproximação final em baixa altitude. O desvio acumulado pelo drone durante o período sem câmera foi de apenas $12{,}53\text{ mm}$, e a aeronave retomou a visão normalmente, concluindo o pouso com um erro final real de apenas $\text{RMSE} = 1{,}75\text{ mm}$ no simulador Gazebo.

Os gráficos de comportamento dinâmico e estatístico do filtro são mostrados na Fig.~\ref{fig:validacao_ekf_4segundos}, destacando o acompanhamento nos eixos horizontais, a resposta contínua de ângulo na volta completa ($S^1$) e a recuperação da trajetória logo após o retorno da imagem.

\begin{figure}[H]
	\centering
	\captionsetup{justification=centering}
	\caption{Desempenho dinâmico do Filtro EKF sob ruído visual, descontinuidade em $S^1$ e oclusão de 4,0 segundos}
	\label{fig:validacao_ekf_4segundos}
	\includegraphics[width=0.85\linewidth]{figuras/validacao_ekf_4segundos.png}
	\vspace{0.2cm}
	\begin{flushleft}
		\footnotesize Fonte: elaboração própria.
	\end{flushleft}
\end{figure}

\subsubsecao{Discussão Crítica sobre o Desempenho do EKF}

A redução de ruído alcançada explica-se pelo bom ajuste do modelo matemático de velocidade constante usado na criação do filtro. Como o drone executa movimentos suaves e amortecidos pelos controladores PD, a aceleração relativa entre a câmera e a armadilha no chão muda pouco a cada intervalo de leitura de $\Delta t = 0{,}05\text{ s}$. Assim, o EKF funciona como um suavizador natural que filtra as oscilações rápidas e aleatórias da detecção sem causar atrasos de resposta que pudessem prejudicar o controle de voo.

Por outro lado, o desvio acumulado durante o bloqueio da câmera ($127{,}55\text{ mm}$ em 4 s) mostra uma limitação natural da previsão baseada apenas no movimento anterior: quando não há imagens da câmera, o erro de estimativa aumenta de acordo com a velocidade em que o drone estava no momento exato em que perdeu o alvo de vista. Se a aeronave sofrer perturbações externas do ar durante esse intervalo sem visão (como rajadas de vento ou turbulência gerada pelas próprias hélices), o cálculo de velocidade constante não consegue perceber essas acelerações imprevistas, acumulando desvios maiores se o bloqueio passar de 4 segundos.

Como oportunidade de melhoria, sugere-se implementar uma fusão de sensores fortemente integrada (\textit{Tightly-Coupled Visual-Inertial EKF}), baseada no método MSCKF de Mourikis e Roumeliotis~\cite{mourikis2007mscfk}, em que as leituras instantâneas dos acelerômetros e giroscópios (unidade IMU) do piloto automático são usadas diretamente na etapa de previsão do filtro de visão. Dessa forma, o sistema conseguiria medir em tempo real qualquer aceleração externa mesmo com a câmera tampada, aumentando a tolerância a perdas de imagem para mais de 6 segundos com um desvio acumulado bem menor.

\subsecao{Teste 3: Retenção no Campo de Visão e Histerese Óptica (Objetivo 4)}

O terceiro teste avaliou o funcionamento da margem de segurança dividida em duas áreas na tela (histerese visual) e a capacidade do sistema de manter a armadilha sempre visível dentro da imagem da câmera durante a descida da Fase 1 ($20{,}0\text{ m} \to 2{,}0\text{ m}$). Para cumprir o Objetivo 4, definiu-se como meta manter o alvo visível $100\%$ do tempo dentro do campo de visão (FOV), sem ultrapassar o limite de segurança de $60\%$ da distância entre o centro e as bordas da tela.

Ao longo das 60 missões autônomas de Monte Carlo realizadas no simulador Gazebo Sim, com o erro inicial de GPS sorteado entre 1,0 m e 4,5 m de distância:
\begin{enumerate}
	\item A taxa de permanência do alvo no campo de visão foi de $100{,}0\%$, sem nenhum caso em que a armadilha tenha saído para fora das bordas da câmera;
	\item A distância média do alvo em relação ao centro da tela foi de $29{,}1\%$ (considerando a metade da largura e da altura da imagem), confirmando que o drone manteve o objeto bem centralizado dentro da área segura ($< 40\%$);
	\item O número de pausas na descida acionadas pela regra de segurança visual (histerese) foi de $0{,}0$ por voo.
\end{enumerate}

\subsubsecao{Discussão Crítica sobre a Histerese e a Dinâmica Descendente}

O fato de não ter ocorrido nenhuma pausa na descida durante os voos normais é resultado direto do bom ajuste da desaceleração vertical aplicada entre 6,0 m e 2,0 m de altura. Conforme definido na Eq.~\ref{eq:fator_frenagem}, a velocidade de descida foi reduzida gradualmente de $1{,}50\text{ m/s}$ para $0{,}15\text{ m/s}$ à medida que o drone se aproximava da altura de troca de fase. Essa redução suave do embalo vertical deu ao controlador horizontal PD tempo suficiente para centralizar o drone sobre a armadilha com folga, mantendo o desvio na tela na faixa média de $29{,}1\%$, sem nunca cruzar o limite de fuga de $60\%$.

O fato de a regra de histerese não precisar ser acionada em voos normais confirma que ela funciona como uma proteção preventiva, ou seja, em condições estáveis, ela não atrasa a descida do drone, porém, se uma rajada de vento lateral ou um erro inicial muito grande empurrasse o alvo para além da margem de $60\%$, a pausa temporária da velocidade vertical $v_z$ evitaria que a câmera perdesse a armadilha de vista até que o drone voltasse a se centralizar.

Para cenários futuros com ventos e turbulências mais fortes, sugere-se como melhoria implementar uma histerese dinâmica adaptativa, ampliando as técnicas de controle com limite de campo de visão estudadas por Thuilot \textit{et al.}~\cite{thuilot2002}. Nesse modelo aprimorado, os limites da área segura e da borda de fuga se ajustariam automaticamente conforme a velocidade horizontal do drone e os ângulos de inclinação para frente, para trás e para os lados (\textit{pitch/roll}) da estrutura, aumentando a margem de segurança durante manobras mais rápidas.

\subsecao{Teste 4: Controle em Cascata e Alinhamento Autônomo de Guinada (Objetivo 5)}

O quarto teste avaliou a estabilidade do sistema de controle que opera em diferentes velocidades de resposta (400 Hz no piloto automático ArduPilot e 20 Hz no módulo do ROS 2) e a precisão do alinhamento automático de giro horizontal (\textit{Yaw}) na altura de 1,50 m durante a Fase 2. A meta estabelecida no Objetivo 5 exigia que o erro angular final em voo controlado fosse menor que $1{,}0^\circ$ ($|e_\psi| \le 1{,}0^\circ$).

A análise conjunta das 60 missões autônomas de Monte Carlo mostrou os seguintes resultados:
\begin{itemize}
	\item \textbf{Erro angular médio absoluto de guinada:} $|e_\psi| = 0{,}46^\circ$;
	\item \textbf{Desvio padrão:} $\sigma_\psi = \pm 0{,}54^\circ$;
	\item \textbf{Cumprimento do critério de estabilização:} $100{,}0\%$ dos voos mantiveram o erro dentro da margem de $|e_\psi| \le 1{,}00^\circ$ sem interrupções por pelo menos $\Delta t \ge 0{,}30\text{ s}$ (6 ciclos seguidos a 20 Hz).
\end{itemize}

\subsubsecao{Discussão Crítica do Alinhamento Angular}

Alcançar um erro médio menor que meio grau ($0{,}46^\circ$), mesmo com o drone começando cada teste apontado para uma direção inicial aleatória qualquer ($\psi_{\text{inicial}} \sim \mathcal{U}(0^\circ; 360^\circ)$), comprova a precisão do método geométrico baseado no formato triangular das três alças. A decisão de fazer esse ajuste de giro com o drone parado no ar (voo pairado) a 1,50 m de altura foi feita já que ao pausar os deslocamentos laterais e a descida durante a rotação, separou-se a força de giro das forças de sustentação vertical, permitindo que o controle interno de estabilidade do ArduPilot amortecesse qualquer balanço restante. Vale ressaltar que a escolha da altitude de 1,50 m para ajustar o \textit{Yaw} se deu devido ao fato de a rede neural do YOLO detectar melhor as alças a essa altura por conta do \textit{dataset} utilizado.

Além disso, o cálculo do giro pelo caminho angular mais curto no círculo completo ($S^1$), realizado pelo filtro \textit{FiltroAnguloEKF}, evitou saltos bruscos nos números dos ângulos. Nos voos em que a frente do drone começava apontada perto de $-180^\circ$ e o ângulo desejado estava em $+93{,}37^\circ$, o algoritmo girou pelo trajeto mais curto, evitando voltas completas desnecessárias de $360^\circ$.

Como ponto de melhoria futura, observa-se que esse método geométrico depende da identificação simultânea das três alças triangulares. Em situações reais em campo, caso a estrutura da armadilha sofra algum amassado ou tenha sujeira cobrindo uma das alças, a marcação das pontas do triângulo pode apresentar falhas. Para contornar esse problema, sugere-se testar modelos de inteligência artificial que estimam diretamente a posição e a inclinação 3D do objeto (pose 6D) por meio de pontos-chave visuais (\textit{Keypoint Detection}). Essa técnica já foi aplicada com sucesso por Ubellacker \textit{et al.}~\cite{ubellacker2024} em sistemas de captura aérea com câmera embarcada, permitindo calcular a orientação espacial do alvo a partir da aparência geral de toda a sua superfície.

\subsecao{Teste 5: Bateria de Monte Carlo da Missão Completa e Acoplamento (Objetivos 6 e 7)}

O quinto teste avaliou o desempenho completo do sistema, do início ao fim da operação, por meio de uma bateria estatística de Monte Carlo com $N = 60$ missões autônomas completas no simulador físico Gazebo Sim conectado ao piloto automático ArduPilot SITL. Cada voo começou com variações aleatórias na altura inicial de cruzeiro ($16{,}0\text{ m a } 20{,}0\text{ m}$), na margem de erro radial do GPS ($1{,}0\text{ m a } 4{,}5\text{ m}$) e na direção inicial para onde o drone estava apontado ($0{,}0^\circ\text{ a } 360{,}0^\circ$).

A Tabela~\ref{tab:consolidado_monte_carlo} reúne os principais resultados de precisão física e tempo de voo extraídos do relatório geral e dos registros de telemetria.

\begin{table}[H]
	\centering
	\footnotesize
	\setlength{\tabcolsep}{4pt}
	\captionsetup{justification=centering}
	\caption{Consolidação estatística das baterias de voo autônomo (Monte Carlo) no simulador}
	\label{tab:consolidado_monte_carlo}
	\vspace{0.2cm}
	\begin{tabular}{>{\raggedright\arraybackslash}p{6.0cm}ccc}
		\hline
		\textbf{Métrica de Desempenho} & \textbf{Valor Obtido} & \textbf{Meta Estabelecida} & \textbf{Status} \\
		\hline
		Total de missões autônomas executadas & 60 & $\ge 20$ voos & Concluído \\
		Taxa de retenção no FOV (Teste 3) & 100,0\% & 100,0\% & Atendido \\
		Afastamento médio da borda útil (FOV) & 29,1\% & $< 60{,}0\%$ (fuga) & Atendido \\
		Bloqueios médios de histerese & 0,0 por voo & N/A & Atendido \\
		Erro médio absoluto de \textit{Yaw} (Teste 4) & $0{,}46^\circ \pm 0{,}54^\circ$ & $\le 1{,}00^\circ$ & Atendido \\
		Erro médio pré-\textit{Blind Drop} (visão em 0,20 m) & $3{,}13 \pm 1{,}12\text{ mm}$ & $\le 5{,}00\text{ mm}$ & Atendido \\
		\textbf{RMSE médio terminal pós-pouso (Gazebo)} & $\mathbf{5{,}25\text{ mm}}$ & $\mathbf{\le 8{,}00\text{ mm}}$ & \textbf{Atendido} \\
		Faixa de dispersão do erro terminal (RMSE) & $[0{,}65; 11{,}04]\text{ mm}$ & $\le 8{,}00\text{ mm}$ & Avaliar \\
		\textbf{Taxa Real de Acoplamento Mecânico} & \textbf{76,7\%} (46/60) & $\mathbf{> 60{,}0\%}$ & \textbf{Superado} \\
		Tempo médio total de missão & $111{,}1\text{ s}$ & $< 90{,}0\text{ s}$ & Avaliar \\
		\hline
	\end{tabular}
	\vspace{0.2cm}
	\begin{flushleft}
		\footnotesize Fonte: elaboração própria.
	\end{flushleft}
\end{table}

A taxa real de encaixe bem-sucedido no final da missão chegou a $76{,}7\%$ (46 capturas concluídas dentro de todos os critérios rigorosos em 60 voos), superando com folga a meta mínima de $> 60{,}0\%$ definida no Objetivo 7. O erro radial médio medido diretamente pelas coordenadas exatas do mundo simulado (\textit{ground truth} físico do Gazebo Sim entre o corpo do drone e o centro da armadilha depois do pouso e do desligamento dos motores) ficou em $5{,}25\text{ mm}$, atendendo com boa margem de segurança ao limite exigido no Objetivo 6 ($\text{RMSE} \le 8{,}00\text{ mm}$).

A distribuição dos erros finais de posição $(e_x, e_y)$ ao redor do centro da armadilha é mostrada no gráfico de dispersão da Fig.~\ref{fig:dispersao_pousos_monte_carlo}, destacando como a maioria dos pousos ficou concentrada dentro do círculo de tolerância máxima ($\le 8{,}0\text{ mm}$).

\begin{figure}[H]
	\centering
	\captionsetup{justification=centering}
	\caption{Diagrama de dispersão métrica dos pontos terminais de pouso nos 60 ensaios de Monte Carlo}
	\label{fig:dispersao_pousos_monte_carlo}
	\includegraphics[width=0.85\linewidth]{figuras/dispersao_pousos_monte_carlo.png}
	\vspace{0.2cm}
	\begin{flushleft}
		\footnotesize Fonte: elaboração própria.
	\end{flushleft}
\end{figure}

\subsubsecao{Discussão Crítica sobre a Discrepância de Erro e o Efeito Solo}

Um ponto importante observado na análise dos resultados foi a diferença entre o erro medido no ar, logo antes da descida final sem câmera, e o erro real medido no chão depois do pouso:
\begin{itemize}
	\item \textbf{Erro pré-\textit{Blind Drop} (em voo a 0,20 m de altitude):} média de $3{,}13\text{ mm} \pm 1{,}12\text{ mm}$;
	\item \textbf{Erro terminal pós-pouso (no solo após modo \texttt{LAND}):} média de $5{,}25\text{ mm}$, com alguns desvios chegando a $11{,}04\text{ mm}$ nos 14 voos que ficaram no limite ($23{,}3\%$).
\end{itemize}

Esse aumento de cerca de 2 milímetros no erro acontece por causa de um fenômeno aerodinâmico chamado Efeito Solo (\textit{Ground Effect}), bastante estudado em drones multirrotores por Sanchez-Cuevas, Heredia e Ollero~\cite{sanchezcuevas2017ground}. Quando o drone coaxial \textit{OctaQuad} chega a menos de 25 cm do topo rígido da armadilha e do chão, o vento empurrado para baixo pelas oito hélices bate contra a superfície e volta para cima formando redemoinhos turbulentos. Essa turbulência empurra o corpo do drone para os lados de forma imprevisível.

A manobra de descida vertical firme a $0{,}80\text{ m/s}$ (\textit{Blind Drop}), mantendo a posição horizontal travada no sistema fixo do terreno (\texttt{LOCAL\_NED}), foi criada justamente para reduzir esse problema, usando o próprio peso e embalo da descida para atravessar essa camada de turbulência até tocar o alvo. Mesmo assim, após encostar no chão demora um pouco até que o piloto automático detecte que o drone está no solo e mande ele desarmar e nesse tempo acontecem pequenos deslizamentos mecânicos de 1 a 2 milímetros já que as hélices continuam girando.

Nos 14 voos em que o erro final ficou entre $8{,}1\text{ mm}$ e $11{,}0\text{ mm}$, a base do mecanismo de captura encostou na borda externa das alças. E como a simulação utilizou um critério geométrico muito rigoroso ($\le 8{,}0\text{ mm}$) para representar o mecanismo sem movimento ativo, esses 14 casos foram identificados como falhas e contabilizados nos testes.

O resultado físico dessa etapa final no ambiente simulado é mostrado na Fig.~\ref{fig:acoplamento_terminal_gazebo}, que exibe o drone pousado sobre a armadilha com o mecanismo de captura encaixado na alça.

\begin{figure}[H]
	\centering
	\captionsetup{justification=centering}
	\caption{Registro do acoplamento terminal entre a interface ventral do VANT e a alça da armadilha no simulador}
	\label{fig:acoplamento_terminal_gazebo}
	\vspace{0.2cm}
	\includegraphics[width=0.85\linewidth]{figuras/acoplamento_terminal_gazebo.png}
	\vspace{0.2cm}
	\begin{flushleft}
		\footnotesize Fonte: elaboração própria.
	\end{flushleft}
\end{figure}

\subsubsecao{Discussão sobre o Tempo Médio de Missão e o Compromisso de Engenharia}

O tempo médio medido para concluir a missão inteira foi de $111{,}1\text{ segundos}$, um pouco acima da previsão inicial de 90 segundos. Essa diferença de cerca de 21 segundos representa uma escolha consciente de projeto (\textit{trade-off}) que prioriza a segurança e a precisão no lugar da pressa:
\begin{enumerate}
	\item Na Fase 1, a rampa de desaceleração entre 6,0 m e 2,0 m ajudou na suavidade da frenagem mas exigiu mais tempo.
	\item Na Fase 2, a parada no ar (\textit{hover}) a 1,50 m para garantir o alinhamento firme de giro ($|e_\psi| \le 1{,}0^\circ$ por pelo menos $0{,}30\text{ s}$) também consumiu bastante tempo para garantir a angulação correta;
	\item Na Fase 3, a etapa de \textbf{ancoragem milimétrica em voo pairado a 0,20 m} exigiu que o drone mantivesse o erro horizontal dentro de $\|e_{xy}\| \le 5\text{ mm}$ durante $2{,}50\text{ segundos}$ seguidos sem interrupções (50 quadros consecutivos a 20 Hz). Se uma pequena oscilação de $5{,}2\text{ mm}$ acontecesse no 48º quadro, o contador de tempo zerava e recomeçava do início como medida de segurança.
\end{enumerate}

Essa exigência rigorosa de estabilidade aumentou a duração total do voo, mas foi exatamente o que garantiu que o drone só iniciasse a descida final quando a chance de erro fosse mínima. Em operações reais de recolhimento, evitar colisões e garantir um encaixe seguro é muito mais importante do que economizar alguns segundos, tornando esse tempo extra totalmente aceitável na prática.

\subsubsecao{Fatores de Melhoria para a Taxa de Acoplamento}

Para aumentar a taxa de sucesso nas capturas de $76{,}7\%$ para mais de $90\%$, recomendam-se duas melhorias complementares:
\begin{enumerate}
	\item \textbf{Aprimoramento Mecânico (Aumento da Abertura e Geometria Autocentrante):} Além da construção física da garra com servomotor ativo e pontas afuniladas ou flexíveis (\textit{soft gripper}), conforme discutido nas revisões sobre manipulação aérea de Bonyan Khamseh, Janabi-Sharifi e Abdessameud~\cite{bonyan2018} e comprovado na prática por Ubellacker \textit{et al.}~\cite{ubellacker2024}, recomenda-se aumentar a abertura física das mandíbulas da garra. No projeto atual, a folga mecânica aceita desvios de até cerca de 8 a 10 mm; ampliar esse vão de entrada aumentará a margem de tolerância do sistema justamente durante a descida final sem câmera (\textit{Blind Drop}), garantindo que os pequenos desvios laterais de até $20\text{ a } 25\text{ mm}$ causados pelo efeito solo continuem dentro da área de captura e sejam guiados mecanicamente para o centro da trava no momento do fechamento;
	
	\item \textbf{Controle Avançado Antidistúrbio:} A troca do controlador PD no trecho final a 0,20 m por algoritmos de controle projetados especificamente para compensar perturbações externas, como o Controle por Modos Deslizantes (SMC), usado por Carvajal \textit{et al.}~\cite{carvajal2024} em sistemas visuais de drones, ou o Controle Preditivo Baseado em Modelo (MPC) aplicado a acoplamentos autônomos~\cite{pramono2025}, garantindo reação rápida contra os empurrões do efeito solo nos últimos 250 milissegundos antes do toque.
\end{enumerate}

\subsecao{Síntese Global e Validação dos Objetivos Específicos}

A Tabela~\ref{tab:sintese_geral_objetivos} reúne o resumo final do cumprimento de todos os sete objetivos específicos definidos na Introdução deste trabalho, comparando as metas previstas com os resultados numéricos obtidos nos experimentos.

\begin{table}[H]
	\centering
	\footnotesize
	\setlength{\tabcolsep}{3pt}
	\captionsetup{justification=centering}
	\caption{Quadro de consolidação final e atendimento aos objetivos específicos da pesquisa}
	\label{tab:sintese_geral_objetivos}
	\vspace{0.2cm}
	\begin{tabular}{>{\raggedright\arraybackslash}p{3.8cm}>{\raggedright\arraybackslash}p{3.4cm}>{\raggedright\arraybackslash}p{3.0cm}>{\raggedright\arraybackslash}p{2.8cm}c}
		\hline
		\textbf{Objetivo Específico} & \textbf{Critério de Projeto} & \textbf{Meta Estipulada} & \textbf{Resultado Alcançado} & \textbf{Status} \\
		\hline
		1. Acurácia dos modelos YOLO & $mAP@50$ nos 3 regimes & $> 90{,}0\%$ em todas as fases & $> 99{,}2\%$ em todas as 3 fases & Superado \\
		2. Latência e taxa de quadros & Inferência unitária (\textit{batch} = 1) & $> 20\text{ fps}$ ($< 50\text{ ms}$) & 28,6 a 51,8 fps ($19{,}3$ a $35\text{ ms}$) & Superado \\
		3. Filtragem EKF e oclusão & Redução ruído e deriva 4 s & Redução $\ge 40\%$; deriva $\le 250$ mm & Red. $51{,}98\%$; deriva $127{,}55$ mm & Superado \\
		4. Retenção no campo de visão & Histerese \textit{dual-zone} 3D & 100,0\% de retenção no FOV & 100,0\% retido (afast. 29,1\%) & Atendido \\
		5. Alinhamento de \textit{Yaw} & Controle em cascata na Fase 2 & $|e_\psi| \le 1{,}00^\circ$ & $|e_\psi| = 0{,}46^\circ \pm 0{,}54^\circ$ & Superado \\
		6. Precisão métrica de captura & RMSE terminal no Gazebo & $\text{RMSE} \le 8{,}00\text{ mm}$ & $\text{RMSE} = 5{,}25\text{ mm}$ (3,13 mm pré-\textit{drop}) & Atendido \\
		7. Taxa de sucesso integrada & Bateria de Monte Carlo ($N=60$) & Sucesso global $> 60{,}0\%$ & \textbf{76,7\%} (46/60 acoplamentos) & Superado \\
		\hline
	\end{tabular}
	\vspace{0.2cm}
	\begin{flushleft}
		\footnotesize Fonte: elaboração própria.
	\end{flushleft}
\end{table}

\subsubsecao{Confronto com a Literatura Correlata}

Ao comparar os resultados obtidos com os estudos apresentados no Capítulo~2, destacam-se avanços práticos importantes:
\begin{itemize}
	\item Em comparação com Bari\v{s}i\'c \textit{et al.}~\cite{barisic2019}, que usaram a rede YOLOv3-Tiny com um filtro linear para perseguir outro drone e alcançaram 60\% de redução de ruído a 20 fps, este trabalho atingiu velocidades de até 51,8 fps com arquiteturas mais modernas (YOLO26n e YOLOv10n). Além disso, o sistema desenvolvido manteve a previsão da trajetória mesmo durante bloqueios de visão de 4,0 segundos e evitou saltos bruscos nos ângulos na transição completa de volta ($S^1$);
	\item Em relação aos estudos clássicos de aproximação e pouso autônomo (como Vidal e Casanova~\cite{vidal2007} e Lopez-Franco \textit{et al.}~\cite{lopezfranco2017}), que trabalharam com precisão na casa dos centímetros (erros típicos entre 30 mm e 100 mm), o método proposto neste trabalho alcançou precisão milimétrica ($\text{RMSE} = 5{,}25\text{ mm}$ e erro médio de giro de $0{,}46^\circ$). Esse nível de exatidão torna possível o encaixe mecânico mesmo com folgas muito pequenas na garra;
	\item Ao contrário de soluções que dependem de câmeras externas de captura de movimento em laboratório (MOCAP, como em Thomas \textit{et al.}~\cite{thomas2014}) ou de etiquetas visuais fiduciais coladas no alvo, e avançando sobre as técnicas de reconhecimento direto sem marcadores (como em Bauer \textit{et al.}~\cite{bauer2023}), o sistema desenvolvido funcionou de forma totalmente autônoma sem precisar alterar a estrutura física da armadilha real. Isso comprova que é viável usar redes de visão computacional embarcadas no próprio drone para realizar tarefas de encaixe aéreo com precisão de milímetros.
\end{itemize}